% Known-source correction: replace only ex:ch14-skeleton-current.
% The historical label is retained; the object is a measure, not delta_E.
\begin{example}[Length measure on the skeleton of a Tate curve]
\label{ex:ch14-skeleton-current}
Let $K$ be a complete nontrivially valued non-Archimedean field, let
$q\in K^\times$ satisfy $0<|q|<1$, and let $E$ be the Tate elliptic curve
with
\[
 E^{\An}=\mathbf G_{m,K}^{\An}/q^{\mathbf Z},
 \qquad \ell=-\log|q|=v_K(q)>0.
\]
As in \Cref{ex:trop-tate-curve}, its canonical skeleton is
$\Sigma=\mathbf R/\ell\mathbf Z$.  Write $i:\Sigma\hookrightarrow E^{\An}$
for the inclusion and $\tau:E^{\An}\to\Sigma$ for the retraction.
These are maps of topological spaces; in particular, the pushforward
of a measure along $i$ below is not an analytic pullback of currents.
The split-uniformization description is also recalled in
\cite[paragraphs~19.2.1--19.2.2]{CLD25}.

For the normalization, the point $\eta_s$ on the skeleton of
$\mathbf G_{m,K}^{\An}$ has Gauss seminorm
\[
 \left|\sum_ma_mT^m\right|_{\eta_s}
       =\max_m|a_m|\,\rme^{-ms},\qquad s\in\mathbf R.
\]
Multiplication by $q$ sends $\eta_s$ to $\eta_{s+\ell}$.
Thus the lattice coordinate $s=-\log|T|$ descends locally to $\Sigma$,
with transition maps given by translations by integer multiples of
$\ell$.  The lattice measure $\rd s$ of
\Cref{def:ch11-lattice-measure} consequently descends to the circle,
and a fundamental interval has length $\ell$.

Define the finite positive Radon measure on $E^{\An}$
\[
                         \mu_\Sigma=i_*(\rd s).
\]
Its associated current acts on functions, not on $(1,1)$-forms:
\[
 \langle[\mu_\Sigma],f\rangle
       =\int_\Sigma f(i(s))\,\rd s,
       \qquad f\in A_{E,c}^{0,0}.
\]
It has bidimension $(0,0)$ and hence bidegree $(1,1)$ on the curve.
Indeed, positivity is immediate and
\[
 |\langle[\mu_\Sigma],f\rangle|
       \leq\ell\,\|f\|_\infty.
\]
This gives an order-zero current, equivalently the construction of
\cite[Proposition~7.2.8]{CLD25} and
\Cref{ex:ch14-forms-and-measures}.  Since $E^{\An}$ is compact, the
constant function $1$ is an admissible test and
\[
                    \langle[\mu_\Sigma],1\rangle=\ell.
\]
The identity
$\int f\,\rd\mu_\Sigma=\int(f\circ i\circ\tau)\,\rd\mu_\Sigma$
for continuous $f$ follows simply from $\tau\circ i=\mathrm{id}_\Sigma$.
It is an identity for this supported measure, not a definition of
$\tau^*$ on arbitrary currents.

By contrast, the integration current of the whole curve is
\[
 \langle\delta_E,\alpha\rangle=\int_{E^{\An}}\alpha,
       \qquad \alpha\in A_{E,c}^{1,1}.
\]
It has bidimension $(1,1)$ and bidegree $(0,0)$.
Thus $[\mu_\Sigma]$ and $\delta_E$ belong to different current spaces;
they cannot be identified, and $\langle\delta_E,1\rangle$ is not a
pairing of the required degree.

Nor does every $(1,1)$-form integrate solely on this fixed circle.
Here is an explicit test.  Let
$p:\mathbf G_{m,K}^{\An}\to E^{\An}$ be the quotient map and put
$B=\{|T-1|<1\}$.  On $B$ one has $|T|=1$, whereas every nontrivial
$q^{\mathbf Z}$-translate has a different value of $|T|$.
Consequently $p|_B$ is an isomorphism onto an ordinary open subset of
$E^{\An}$.  This subset is disjoint from $\Sigma$: at a Gauss point
$\eta_s$ one has $|T-1|_{\eta_s}=\max\{\rme^{-s},1\}\geq1$.

Choose $0<r<R<1$ and a nonnegative, nonzero function
$\chi\in C_c^\infty((-\log R,-\log r))$.
On the open annulus $r<|T-1|<R$ write $u=-\log|T-1|$ and take
\[
                           \alpha=\chi(u)\,\dprime u\wedge\ddprime u.
\]
Transport this form by $p|_B$ and extend it by zero to $E^{\An}$.
This is a smooth compactly supported form: its support is contained
in the image of a closed subannulus strictly inside $B$.
The coordinate $T-1$ identifies the annulus with a standard annulus,
of calibration weight $1$, so
\[
 \langle\delta_E,\alpha\rangle
       =\int_{-\log R}^{-\log r}\chi(u)\,\rd u>0,
                  \qquad \alpha|_\Sigma=0.
\]
Thus the circle length $\ell$ is the mass of the specifically defined
measure $\mu_\Sigma$, not a function-test mass for the full
integration current $\delta_E$.
\end{example}
